% =====================================================================
%  ĮVADAS — rašomas paskutinis (rugs. 18), bet failas reikalingas dabar,
%  kad ataskaita.tex kompiliuotųsi.
%  Tikslinė apimtis: 1–2 psl.
% =====================================================================

\subsection*{Temos aktualumas}

% TODO: IoT įrenginių skaičiaus augimas, jų silpna apsauga,
% Mirai tipo botnetų mastas. 2–3 pastraipos su nuorodomis.

\subsection*{Darbo tikslas}

Sukurti ir eksperimentiškai įvertinti dirbtiniu intelektu pagrįstą IoT tinklo
kibernetinių atakų aptikimo sprendimą.

\subsection*{Darbo uždaviniai}

\begin{enumerate}
  \item Išanalizuoti IoT tinklams būdingas kibernetines atakas ir jų aptikimo metodus.
  \item Išnagrinėti dirbtinio intelekto ir mašininio mokymosi metodų taikymo
        galimybes IoT tinklo atakoms aptikti.
  \item Parinkti ir pagrįsti tinkamus DI metodus IoT tinklo anomalijoms ir
        kibernetinėms atakoms identifikuoti.
  \item Sukurti DI pagrįstą IoT tinklo atakų aptikimo sprendimą.
  \item Eksperimentiškai įvertinti sukurto sprendimo tikslumą, klaidingų
        aptikimų skaičių ir kitus veikimo rodiklius.
  \item Palyginti skirtingų DI metodų efektyvumą aptinkant IoT tinklo atakas.
\end{enumerate}

\subsection*{Darbo struktūra}

% TODO (rugs. 18): 1 pastraipa — ką aprašo kiekvienas skyrius.
